% Session 6 -- bipartite order versus geometric frustration.
% Placement: sec. "Bipartite order versus geometric frustration", after the triangular-plaquette
% paragraph that introduces the 120-degree compromise.
%
% Energies are quoted per plaquette for classical spins of length S, so that
% S_i . S_j = S^2 cos(theta_ij): a satisfied (antiparallel) bond contributes -J S^2,
% a parallel bond +J S^2, and a 120-degree bond -J S^2 / 2.
\begin{aqmfigure}{fig:s06:frustration}{How the geometry decides if an antiferromagnet can order.
(a) On a bipartite lattice the sites can be coloured with two colours, every bond joins the two
colours, and the N\'eel state of Section 6.10 satisfies all bonds at the same time. (b) On a triangle
it is not possible to use two colours. If we fix two spins antiparallel, the third one is parallel to
one of its neighbours for any choice, and one bond stays frustrated. For classical spins, this collinear compromise costs $+JS^2$ on the frustrated bond, and the total is
$-JS^2$ per triangle. (c) The energy is lower if no single bond is fully satisfied. Three spins at $120^\circ$ give $\mathbf
S_i\cdot\mathbf S_j=-S^2/2$ on \emph{every} bond and a total of $-\tfrac32 JS^2$, which is lower than
(b) by half of its value. This is the origin of the noncollinear order to which the mean-field
solution of the corresponding electronic Hamiltonian also converges.}

\begin{tikzpicture}[x=1cm, y=1cm]

% --- local helper: a classical spin, dot at (#1,#2), arrow at angle #3 -------
\newcommand{\clspin}[4]{% x, y, angle, colour
  \draw[aqm, draw=#4, -{Stealth[length=4.6pt,width=3.8pt]}]
    ({#1-0.26*cos(#3)},{#2-0.26*sin(#3)}) -- ({#1+0.28*cos(#3)},{#2+0.28*sin(#3)});}

% --- panel headings, all on one baseline ------------------------------------
\node[aqmpanel, anchor=south] at (0.92,2.45) {(a) bipartite};
\node[aqmpanel, anchor=south] at (4.90,2.45) {(b) frustrated};
\node[aqmpanel, anchor=south] at (9.60,2.45) {(c) the $120^\circ$ state};

% ============================ (a) bipartite ==================================
\begin{scope}[shift={(0,0)}]
  \def\dd{0.92}
  \foreach \i in {0,1,2} {
    \draw[aqmbond] (0,\i*\dd) -- (2*\dd,\i*\dd);
    \draw[aqmbond] (\i*\dd,0) -- (\i*\dd,2*\dd);
  }
  % sublattice A: (i+j) even, spins up
  \foreach \i/\j in {0/0,0/2,1/1,2/0,2/2} {
    \node[aqmsite, minimum size=5pt] at (\i*\dd,\j*\dd) {};
    \clspin{\i*\dd}{\j*\dd}{90}{aqmaccent}
  }
  % sublattice B: (i+j) odd, spins down
  \foreach \i/\j in {0/1,1/0,1/2,2/1} {
    \node[aqmsiteopen, minimum size=5pt] at (\i*\dd,\j*\dd) {};
    \clspin{\i*\dd}{\j*\dd}{270}{aqmaccent}
  }
  \node[aqmlabsmall, anchor=north, align=center, text width=4.2cm] at (\dd,-0.55)
    {every bond $\mathbf S_i\cdot\mathbf S_j=-S^2$,\\[1pt]
     \textcolor{aqmaccent}{all satisfied at once}};
\end{scope}

% ============================ (b) frustrated triangle ========================
\begin{scope}[shift={(4.90,0.62)}]
  \def\rr{0.86}
  \coordinate (b1) at (90:\rr);
  \coordinate (b2) at (210:\rr);
  \coordinate (b3) at (330:\rr);
  \draw[aqmbond] (b1) -- (b2);
  \draw[aqmbond] (b2) -- (b3);
  \draw[aqmbondstrong, draw=aqmbroken] (b1) -- (b3);

  \node[aqmsite, minimum size=5pt] at (b1) {};
  \node[aqmsiteopen, minimum size=5pt] at (b2) {};
  \node[aqmsite, minimum size=5pt] at (b3) {};
  \clspin{0}{0.860}{90}{aqmaccent}
  \clspin{-0.745}{-0.430}{270}{aqmaccent}
  \clspin{0.745}{-0.430}{90}{aqmaccent}

  % mark the unsatisfiable bond with a cross at its midpoint
  \draw[aqm, draw=aqmbroken, line width=1.0pt] (0.272,0.115) -- (0.472,0.315);
  \draw[aqm, draw=aqmbroken, line width=1.0pt] (0.272,0.315) -- (0.472,0.115);
  \node[aqmlabsmall, text=aqmbroken, anchor=west] at (0.55,0.30) {$+JS^2$};

  \node[aqmlabsmall, anchor=north, align=center, text width=3.6cm] at (0,-1.17)
    {two bonds $-JS^2$, one bond $+JS^2$,\\[1pt]
     total \textcolor{aqmbroken}{$-JS^2$}};
\end{scope}

% ============================ (c) 120-degree order ===========================
\begin{scope}[shift={(9.60,0.62)}]
  \def\rr{0.86}
  \coordinate (c1) at (90:\rr);
  \coordinate (c2) at (210:\rr);
  \coordinate (c3) at (330:\rr);
  \draw[aqmbond] (c1) -- (c2);
  \draw[aqmbond] (c2) -- (c3);
  \draw[aqmbond] (c1) -- (c3);

  \node[aqmsite, minimum size=5pt] at (c1) {};
  \node[aqmsite, minimum size=5pt] at (c2) {};
  \node[aqmsite, minimum size=5pt] at (c3) {};
  \clspin{0}{0.860}{90}{aqmbroken}
  \clspin{-0.745}{-0.430}{210}{aqmbroken}
  \clspin{0.745}{-0.430}{330}{aqmbroken}

  % the relative angle between two of the three spin directions
  \draw[aqmguide] (0,0) -- (90:0.56);
  \draw[aqmguide] (0,0) -- (330:0.56);
  \draw[aqm, draw=aqmmuted, line width=\aqmthinwidth] (330:0.28) arc (-30:90:0.28);
  \draw[aqmguide] (0.94,0.46) -- (0.27,0.15);
  \node[aqmlabsmall, text=aqmmuted, anchor=west] at (0.97,0.46) {$120^\circ$};

  \node[aqmlabsmall, anchor=north, align=center, text width=3.6cm] at (0,-1.17)
    {every bond $-\tfrac12 JS^2$,\\[1pt]
     total \textcolor{aqmbroken}{$-\tfrac32 JS^2$}: lower};
\end{scope}

\end{tikzpicture}

\end{aqmfigure}
